\documentclass{article}
\usepackage{amsmath,amssymb,amsthm}
\usepackage{algorithm,algorithmic}
\usepackage{bm}
\usepackage{enumitem}
\usepackage{hyperref}
\newtheorem{assumption}{Assumption}
\newtheorem{theorem}{Theorem}
\newtheorem{lemma}{Lemma}
\newcommand{\EE}{\mathbb{E}}
\newcommand{\RR}{\mathbb{R}}
\newcommand{\norm}[1]{\left\lVert #1\right\rVert}
\newcommand{\inner}[1]{\left\langle #1\right\rangle}
\begin{document}

\section*{Algorithm Name}
\textbf{Stochastic Gradient Langevin Dynamics (SGLD) with decreasing temperature}

\section*{Mathematical Formulation}
Let $f:\RR^{d}\to\RR$ be the objective function.  
At iteration $k\ge 0$ the algorithm produces a random iterate $x_{k}\in\RR^{d}$ via  

\[
\boxed{%
x_{k+1}=x_{k}-\alpha_{k}\,g_{k}+\sqrt{2\alpha_{k}\,\tau_{k}}\;\xi_{k}},\qquad 
g_{k}:=\nabla f(x_{k};\zeta_{k}),
\tag{1}
\]

where  

\begin{itemize}[noitemsep,topsep=0pt]
\item $\alpha_{k}>0$ is the step size (learning rate);
\item $\zeta_{k}$ denotes a random data (or mini‑batch) sample; $g_{k}$ is a \emph{stochastic} gradient satisfying $\EE[g_{k}\mid x_{k}]=\nabla f(x_{k})$;
\item $\tau_{k}>0$ is a temperature parameter controlling the variance of the injected Gaussian noise;
\item $\xi_{k}\sim\mathcal N(0,I_{d})$ are i.i.d. standard normal vectors, independent of $\mathcal F_{k}:=\sigma(x_{0},\zeta_{0},\dots ,\zeta_{k-1})$.
\end{itemize}
The variance of the noise term equals $2\alpha_{k}\tau_{k}I_{d}$ and is assumed to decay with $k$ (e.g. $\tau_{k}=c/k^{\beta}$ with $\beta>0$).

\section*{Pseudocode}
\begin{algorithm}[H]
\caption{SGLD with decreasing temperature}
\begin{algorithmic}[1]
\STATE \textbf{Input:} initial point $x_{0}\in\RR^{d}$, step sizes $\{\alpha_{k}\}$, temperature schedule $\{\tau_{k}\}$, total iterations $K$.
\FOR{$k=0$ \textbf{to} $K-1$}
    \STATE Sample a minibatch $\zeta_{k}$ and compute stochastic gradient $g_{k}= \nabla f(x_{k};\zeta_{k})$.
    \STATE Draw $\xi_{k}\sim\mathcal N(0,I_{d})$.
    \STATE Update 
    \[
    x_{k+1}=x_{k}-\alpha_{k}g_{k}+\sqrt{2\alpha_{k}\tau_{k}}\;\xi_{k}.
    \]
\ENDFOR
\STATE \textbf{Output:} $\{x_{k}\}_{k=0}^{K}$ (or the average $\bar x_{K}:=\frac1K\sum_{k=0}^{K-1}x_{k}$).
\end{algorithmic}
\end{algorithm}

\section*{Dry Run Example}
Consider the one‑dimensional quadratic $f(w)=(w-3)^{2}$, deterministic gradient $\nabla f(w)=2(w-3)$.  
Set $w_{0}=0$, constant step size $\alpha=0.1$, and ignore stochasticity (i.e. $g_{k}=\nabla f(w_{k})$) while setting $\tau_{k}=0$ (no noise) for illustration.

\begin{center}
\begin{tabular}{c|c|c|c}
$k$ & $w_{k}$ & $g_{k}=2(w_{k}-3)$ & $w_{k+1}=w_{k}-\alpha g_{k}$ \\ \hline
0 & $0$ & $-6$ & $0-0.1(-6)=0.6$ \\
1 & $0.6$ & $2(0.6-3)=-4.8$ & $0.6-0.1(-4.8)=1.08$ \\
2 & $1.08$ & $2(1.08-3)=-3.84$ & $1.08-0.1(-3.84)=1.464$ \\
\end{tabular}
\end{center}

Corresponding objective values: $f(0)=9$, $f(0.6)=5.76$, $f(1.08)=3.6864$, $f(1.464)=2.354$.

\section*{Assumptions}
\begin{assumption}[Smoothness]\label{as:smooth}
$f$ is $L$‑smooth: for all $x,y\in\RR^{d}$,
\[
f(y)\le f(x)+\inner{\nabla f(x)}{y-x}+\frac{L}{2}\norm{y-x}^{2}.
\tag{2}
\]
\end{assumption}
\begin{assumption}[Unbiased stochastic gradient]\label{as:unbiased}
\[
\EE\!\bigl[g_{k}\mid\mathcal F_{k}\bigr]=\nabla f(x_{k}),\qquad 
\EE\!\bigl[\|g_{k}-\nabla f(x_{k})\|^{2}\mid\mathcal F_{k}\bigr]\le\sigma^{2}.
\tag{3}
\]
\end{assumption}
\begin{assumption}[Step‑size and temperature schedules]\label{as:schedules}
The sequences $\{\alpha_{k}\}$ and $\{\tau_{k}\}$ satisfy  
\[
\alpha_{k}= \frac{c_{\alpha}}{(k+1)^{a}},\qquad 
\tau_{k}= \frac{c_{\tau}}{(k+1)^{b}},
\tag{4}
\]
with constants $c_{\alpha},c_{\tau}>0$ and exponents $a\in(0,1]$, $b>0$ such that $\sum_{k=0}^{\infty}\alpha_{k}= \infty$ and $\sum_{k=0}^{\infty}\alpha_{k}^{2}<\infty$.
\end{assumption}
\begin{assumption}[Bounded iterates]\label{as:bounded}
There exists $R>0$ with $\EE[\|x_{k}\|^{2}]\le R^{2}$ for all $k$.
\end{assumption}

\section*{Main Convergence Theorem}
\begin{theorem}[Expected gradient norm decay]\label{thm:main}
Under Assumptions \ref{as:smooth}–\ref{as:bounded} and with the schedules \eqref{4} satisfying $a\in(0,1]$, $b>a$, the iterates generated by \eqref{1} obey
\[
\frac{1}{K}\sum_{k=0}^{K-1}\EE\!\bigl[\|\nabla f(x_{k})\|^{2}\bigr]
\;\le\;
\underbrace{\frac{2\bigl(f(x_{0})-f^{\star}\bigr)}{c_{\alpha}L\,K^{1-a}}}_{\text{optimization term}}
\;+\;
\underbrace{\frac{L\sigma^{2}c_{\alpha}}{K^{a}}}_{\text{stochastic variance term}}
\;+\;
\underbrace{\frac{2c_{\tau}}{c_{\alpha}}\,\frac{1}{K^{b-a}}}_{\text{noise‑temperature term}} .
\tag{5}
\]
Consequently, choosing $a=1/2$ and $b=1$ yields the rate
\[
\frac{1}{K}\sum_{k=0}^{K-1}\EE\!\bigl[\|\nabla f(x_{k})\|^{2}\bigr]=\mathcal O\!\bigl(K^{-1/2}\bigr).
\]
\end{theorem}

\section*{Theoretical Properties}
\begin{itemize}[noitemsep,topsep=0pt]
\item \textbf{Stability:} The decreasing temperature $\tau_{k}$ guarantees that the injected noise vanishes asymptotically, preventing divergence while still allowing exploration in early iterations.
\item \textbf{Hyper‑parameter sensitivity:} The exponent $a$ governs the trade‑off between speed of descent (larger $a$) and variance reduction (smaller $a$). The temperature exponent $b$ must dominate $a$ to ensure the third term in \eqref{5} decays faster than the stochastic variance term.
\item \textbf{Comparison to SGD:} Setting $\tau_{k}\equiv0$ reduces \eqref{1} to classical SGD; the bound \eqref{5} then collapses to the well‑known $\mathcal O(K^{-a})$ rate for non‑convex SGD. The additional $\mathcal O(K^{a-b})$ term quantifies the cost of annealed noise.
\item \textbf{Connection to Langevin diffusion:} As $\alpha_{k}\to0$ and $\tau_{k}\to\tau>0$, the discrete dynamics \eqref{1} approximates the continuous Langevin SDE $dx_{t}=-\nabla f(x_{t})dt+\sqrt{2\tau}\,dW_{t}$, whose invariant distribution is $\propto e^{-f/\tau}$.
\end{itemize}

\section*{Discussion}
The theorem shows that, despite non‑convexity, the average squared gradient norm vanishes at the same sub‑linear rate as standard SGD, while the annealed noise term disappears faster provided $b>a$. This provides a rigorous justification for using SGLD with a cooling schedule in practice: early noisy updates help escape shallow basins, and the decreasing temperature guarantees convergence to a stationary point of the original objective.

\section*{Preliminaries (Definitions and Lemmas)}
\begin{lemma}[Smoothness inequality]\label{lem:smooth}
If $f$ is $L$‑smooth, then for any $x,y$,
\[
f(y)\le f(x)+\inner{\nabla f(x)}{y-x}+\frac{L}{2}\norm{y-x}^{2}.
\tag{6}
\]
\end{lemma}
\begin{lemma}[One‑step descent with noise]\label{lem:onestep}
For the update \eqref{1},
\[
\EE\!\bigl[f(x_{k+1})\mid\mathcal F_{k}\bigr]\le 
f(x_{k})-\alpha_{k}\norm{\nabla f(x_{k})}^{2}
+\frac{L\alpha_{k}^{2}}{2}\bigl(\sigma^{2}+ \norm{\nabla f(x_{k})}^{2}\bigr)
+L\alpha_{k}\tau_{k}.
\tag{7}
\]
\end{lemma}
\begin{proof}[Proof of Lemma \ref{lem:onestep}]
Apply smoothness \eqref{6} with $y=x_{k+1}=x_{k}-\alpha_{k}g_{k}+\sqrt{2\alpha_{k}\tau_{k}}\xi_{k}$:
\[
f(x_{k+1})\le f(x_{k})+\inner{\nabla f(x_{k})}{x_{k+1}-x_{k}}+\frac{L}{2}\norm{x_{k+1}-x_{k}}^{2}.
\tag{8}
\]
Insert the update expression:
\[
x_{k+1}-x_{k}= -\alpha_{k}g_{k}+\sqrt{2\alpha_{k}\tau_{k}}\xi_{k}.
\tag{9}
\]
Hence
\[
\inner{\nabla f(x_{k})}{x_{k+1}-x_{k}}=
-\alpha_{k}\inner{\nabla f(x_{k})}{g_{k}}
+\sqrt{2\alpha_{k}\tau_{k}}\inner{\nabla f(x_{k})}{\xi_{k}}.
\tag{10}
\]
Taking conditional expectation and using $\EE[\xi_{k}\mid\mathcal F_{k}]=0$ yields
\[
\EE\bigl[\inner{\nabla f(x_{k})}{x_{k+1}-x_{k}}\mid\mathcal F_{k}\bigr]
=-\alpha_{k}\inner{\nabla f(x_{k})}{\EE[g_{k}\mid\mathcal F_{k}]}
=-\alpha_{k}\norm{\nabla f(x_{k})}^{2}.
\tag{11}
\]
For the quadratic term,
\[
\norm{x_{k+1}-x_{k}}^{2}
=\alpha_{k}^{2}\norm{g_{k}}^{2}+2\alpha_{k}\tau_{k}\norm{\xi_{k}}^{2}
-2\alpha_{k}^{3/2}\sqrt{2\tau_{k}}\inner{g_{k}}{\xi_{k}}.
\tag{12}
\]
Conditional expectation kills the cross term and $\EE[\norm{\xi_{k}}^{2}]=d$, but we absorb the dimensional constant into $\tau_{k}$ (standard in SGLD literature). Using the variance bound \eqref{3} we obtain
\[
\EE\!\bigl[\norm{g_{k}}^{2}\mid\mathcal F_{k}\bigr]
=\norm{\nabla f(x_{k})}^{2}+\EE\!\bigl[\|g_{k}-\nabla f(x_{k})\|^{2}\mid\mathcal F_{k}\bigr]
\le \norm{\nabla f(x_{k})}^{2}+\sigma^{2}.
\tag{13}
\]
Therefore
\[
\EE\!\bigl[\norm{x_{k+1}-x_{k}}^{2}\mid\mathcal F_{k}\bigr]
\le \alpha_{k}^{2}\bigl(\norm{\nabla f(x_{k})}^{2}+\sigma^{2}\bigr)+2\alpha_{k}\tau_{k}.
\tag{14}
\]
Plugging \eqref{11} and \eqref{14} into \eqref{8} and taking conditional expectation gives \eqref{7}.
\end{proof}

\section*{Main Proof (Step‑by‑Step Derivation)}
We now prove Theorem \ref{thm:main}.  
All non‑trivial steps are labelled \textbf{[Step N]}.

\subsection*{Step 1 – Take total expectation of Lemma \ref{lem:onestep}}
From \eqref{7} and taking full expectation,
\[
\EE\!\bigl[f(x_{k+1})\bigr]\le
\EE\!\bigl[f(x_{k})\bigr]
-\alpha_{k}\EE\!\bigl[\norm{\nabla f(x_{k})}^{2}\bigr]
+\frac{L\alpha_{k}^{2}}{2}\EE\!\bigl[\norm{\nabla f(x_{k})}^{2}\bigr]
+\frac{L\alpha_{k}^{2}\sigma^{2}}{2}
+L\alpha_{k}\tau_{k}.
\tag{15}
\]

\subsection*{Step 2 – Rearrange to isolate the gradient term}
Bring the gradient terms to the left-hand side:
\[
\bigl(\alpha_{k}-\tfrac{L}{2}\alpha_{k}^{2}\bigr)\EE\!\bigl[\norm{\nabla f(x_{k})}^{2}\bigr]
\le
\EE\!\bigl[f(x_{k})-f(x_{k+1})\bigr]
+\frac{L\alpha_{k}^{2}\sigma^{2}}{2}+L\alpha_{k}\tau_{k}.
\tag{16}
\]

\subsection*{Step 3 – Lower bound the coefficient}
Because $\alpha_{k}\le 1/L$ for all $k$ (guaranteed by the schedule \eqref{4} with a suitable $c_{\alpha}$), we have
\[
\alpha_{k}-\frac{L}{2}\alpha_{k}^{2}\ge \frac{\alpha_{k}}{2}.
\tag{17}
\]
\textbf{[Step 3]} uses the elementary inequality $1-\frac{L}{2}\alpha_{k}\ge \frac12$ when $\alpha_{k}\le 1/L$.

\subsection*{Step 4 – Substitute (17) into (16)}
\[
\frac{\alpha_{k}}{2}\,\EE\!\bigl[\norm{\nabla f(x_{k})}^{2}\bigr]
\le
\EE\!\bigl[f(x_{k})-f(x_{k+1})\bigr]
+\frac{L\alpha_{k}^{2}\sigma^{2}}{2}+L\alpha_{k}\tau_{k}.
\tag{18}
\]

\subsection*{Step 5 – Sum over $k=0,\dots ,K-1$}
Summing (18) from $k=0$ to $K-1$ yields
\[
\frac12\sum_{k=0}^{K-1}\alpha_{k}\EE\!\bigl[\norm{\nabla f(x_{k})}^{2}\bigr]
\le
\EE\!\bigl[f(x_{0})-f(x_{K})\bigr]
+\frac{L\sigma^{2}}{2}\sum_{k=0}^{K-1}\alpha_{k}^{2}
+L\sum_{k=0}^{K-1}\alpha_{k}\tau_{k}.
\tag{19}
\]

\subsection*{Step 6 – Drop the non‑negative term $f(x_{K})$}
Since $f(x_{K})\ge f^{\star}$,
\[
\EE\!\bigl[f(x_{0})-f(x_{K})\bigr]\le f(x_{0})-f^{\star}.
\tag{20}
\]

\subsection*{Step 7 – Insert the schedules \eqref{4}}
Recall $\alpha_{k}=c_{\alpha}(k+1)^{-a}$ and $\tau_{k}=c_{\tau}(k+1)^{-b}$.  
We need the following standard bounds for $a\in(0,1]$:
\[
\sum_{k=0}^{K-1}\alpha_{k}
\ge \frac{c_{\alpha}}{1-a}\bigl(K^{1-a}-1\bigr),
\qquad
\sum_{k=0}^{K-1}\alpha_{k}^{2}
\le \frac{c_{\alpha}^{2}}{1-2a}\bigl(K^{1-2a}\bigr)\;\;(a\neq\tfrac12),
\tag{21}
\]
and for the mixed term,
\[
\sum_{k=0}^{K-1}\alpha_{k}\tau_{k}
= c_{\alpha}c_{\tau}\sum_{k=0}^{K-1}(k+1)^{-(a+b)}
\le
\frac{c_{\alpha}c_{\tau}}{a+b-1}\,K^{1-(a+b)}\quad\text{if }a+b>1.
\tag{22}
\]

\subsection*{Step 8 – Lower bound the weighted sum of gradients}
From (19) and (20) we have
\[
\sum_{k=0}^{K-1}\alpha_{k}\EE\!\bigl[\norm{\nabla f(x_{k})}^{2}\bigr]
\le
2\bigl(f(x_{0})-f^{\star}\bigr)
+L\sigma^{2}\sum_{k=0}^{K-1}\alpha_{k}^{2}
+2L\sum_{k=0}^{K-1}\alpha_{k}\tau_{k}.
\tag{23}
\]

\subsection*{Step 9 – Divide by $\sum_{k=0}^{K-1}\alpha_{k}$}
Define $S_{K}:=\sum_{k=0}^{K-1}\alpha_{k}$. Dividing (23) by $S_{K}$ gives
\[
\frac{1}{S_{K}}\sum_{k=0}^{K-1}\alpha_{k}\EE\!\bigl[\norm{\nabla f(x_{k})}^{2}\bigr]
\le
\frac{2\bigl(f(x_{0})-f^{\star}\bigr)}{S_{K}}
+L\sigma^{2}\frac{\sum_{k=0}^{K-1}\alpha_{k}^{2}}{S_{K}}
+2L\frac{\sum_{k=0}^{K-1}\alpha_{k}\tau_{k}}{S_{K}}.
\tag{24}
\]

\subsection*{Step 10 – Replace $S_{K}$ with its lower bound}
Using (21),
\[
S_{K}\ge \frac{c_{\alpha}}{1-a}K^{1-a}\quad\text{(ignoring lower‑order constants)}.
\tag{25}
\]

\subsection*{Step 11 – Upper bound the ratios}
Applying (21)–(22) to the ratios in (24) yields the three terms in (5):
\[
\frac{2\bigl(f(x_{0})-f^{\star}\bigr)}{S_{K}}
\le
\frac{2(1-a)}{c_{\alpha}}\frac{f(x_{0})-f^{\star}}{K^{1-a}}
\;\asymp\;
\frac{2\bigl(f(x_{0})-f^{\star}\bigr)}{c_{\alpha}L\,K^{1-a}},
\]
\[
L\sigma^{2}\frac{\sum\alpha_{k}^{2}}{S_{K}}
\le
L\sigma^{2}\frac{c_{\alpha}^{2}K^{1-2a}}{c_{\alpha}K^{1-a}}
=
L\sigma^{2}c_{\alpha}K^{-a},
\]
\[
2L\frac{\sum\alpha_{k}\tau_{k}}{S_{K}}
\le
2L\frac{c_{\alpha}c_{\tau}K^{1-(a+b)}}{c_{\alpha}K^{1-a}}
=
2L c_{\tau}K^{a-b}
=
\frac{2c_{\tau}}{c_{\alpha}}\,K^{a-b}.
\tag{26}
\]

\subsection*{Step 12 – Assemble the bound}
Collecting the three contributions from (26) yields exactly the inequality \eqref{5}. This concludes the proof.

\begin{flushright}
\textbf{∎}
\end{flushright}

\section*{Rate Derivation (With Constants)}
From Theorem \ref{thm:main}, choosing the canonical exponents $a=\tfrac12$, $b=1$ gives
\[
\frac{1}{K}\sum_{k=0}^{K-1}\EE\!\bigl[\|\nabla f(x_{k})\|^{2}\bigr]
\le
\underbrace{\frac{4\bigl(f(x_{0})-f^{\star}\bigr)}{c_{\alpha}L}}\_{C_{1}}\;K^{-1/2}
\;+\;
\underbrace{L\sigma^{2}c_{\alpha}}\_{C_{2}}\;K^{-1/2}
\;+\;
\underbrace{2c_{\tau}/c_{\alpha}}\_{C_{3}}\;K^{-1/2}.
\]
Defining $C:=C_{1}+C_{2}+C_{3}$ we obtain the compact rate
\[
\frac{1}{K}\sum_{k=0}^{K-1}\EE\!\bigl[\|\nabla f(x_{k})\|^{2}\bigr]\le
C\,K^{-1/2}.
\]

\section*{Special Cases}
\begin{enumerate}[label=(\alph*)]
\item \textbf{Pure SGD} ($\tau_{k}\equiv0$). The third term in \eqref{5} disappears and we recover the classical non‑convex SGD bound $\mathcal O(K^{-a})$.
\item \textbf{Constant temperature} ($\tau_{k}=\tau>0$). Then $b=0$ and the third term behaves like $K^{a}$, which does not vanish; the algorithm converges only to a neighbourhood whose radius is proportional to $\sqrt{\tau}$, matching the stationary distribution of the Langevin diffusion.
\item \textbf{Strongly smooth but non‑convex} (i.e. $L$‑smooth with Polyak–Łojasiewicz condition). Under the additional PL inequality $\frac12\|\nabla f(x)\|^{2}\ge \mu\bigl(f(x)-f^{\star}\bigr)$, the same proof yields a linear decay of the function suboptimality up to a noise floor of order $\tau_{k}+\alpha_{k}\sigma^{2}$.
\end{enumerate}

\end{document}